\section{Introduction}
\label{sec:introduction}

Text embedding models support retrieval, semantic similarity, clustering, and
classification, but the strongest encoders are often too costly for
latency- or memory-constrained deployment. Knowledge distillation offers a
natural route to compact embedding models by transferring supervision from a
frozen teacher to a smaller student
\citep{hinton2015distilling,reimers2020multilingual}. Unlike class logits,
however, embedding coordinates have no shared semantics across models. Teacher
and student representations may differ in dimension, orientation, and
architecture. Existing methods address this mismatch with learned projectors,
relational or contrastive objectives, and intermediate-layer supervision
\citep{park2019rkd,tian2020crd,xu2023distillcse,zhang2024jasper,truong2025emo,talas}.
Hierarchy-based transfer can be effective, but it requires access to internal
teacher representations and assumptions about cross-model layer
correspondence \citep{ko2023revisiting,kim2025layerselection}. This machinery
can obscure a more basic question: whether teacher and student endpoints have
first been expressed in comparable coordinates. We therefore study the endpoint
interface as a potential bottleneck in cross-dimensional distillation.

We test this hypothesis by centering the transfer mechanism on an
\emph{extrinsic interface} that expresses teacher representations in coordinates
the student can match pointwise. We augment this primary mechanism with an
\emph{intrinsic regularizer} defined on properties of the representation cloud
that do not depend on its coordinates or ambient dimension. This decomposition
leads to \textbf{GATE-KD}, Gauge Alignment and Topology for Embedding
Distillation. GATE-KD fits a fixed principal-component subspace to cached teacher
endpoints to provide a non-trainable reduction at the student width. Within the
retained subspace, orthogonal Procrustes alignment fits the orientation of the
teacher targets to the student's initial coordinates in closed form, without
changing sample correspondence or pairwise geometry.
Because projection cannot preserve all teacher geometry, GATE-KD adds a
degree-zero persistent-homology ($H_0$) summary regularizer that matches
distributions of connectivity scales in the native teacher and student spaces.
This coarse auxiliary summary is invariant to coordinates and ambient dimension;
it does not preserve sample-level neighborhoods or the full metric geometry. The resulting objective
supervises only the final student embedding, trains no alignment module, and
requires no online teacher forward pass during ordinary training.

Experiments across three teacher--student configurations and nine sentence-level
benchmarks under a label-free task-adapted protocol provide evidence consistent
with this endpoint-centered view. A per-minibatch rotation-invariant control
trails GATE-KD by 2.05 points; separate schedule controls suggest that
consistent corpus-level targets matter beyond allowing an adaptive rotation,
while periodic refitting provides only a smaller additional gain. Native-space
$H_0$ adds 0.53 points over endpoint-only training and performs best among the
evaluated structural controls and weights. Overall, GATE-KD obtains the
numerically highest aggregate score among the evaluated distilled students in
all three configurations, exceeding the next-best baseline by 0.38, 0.48, and
1.59 points. Within the measured minibatch optimization loop, it also records
the lowest update time and cumulative optimization time among the evaluated
methods in all three configurations. These results show that
final-embedding-only distillation can be competitive with the evaluated
hierarchy-based alternatives in the task-adapted sentence-level settings
studied here.

We make three contributions.
\begin{itemize}
    \item We formulate cross-dimensional embedding distillation around a
    comparable final-embedding interface, separating pointwise coordinate
    comparability from auxiliary coordinate-invariant structure.
    \item We introduce a fixed-reduction, closed-form gauge interface using only
    cached teacher endpoints, together with native-space $H_0$ regularization as
    a complementary signal and epoch-wise gauge refresh as a secondary
    optimization refinement.
    \item We provide a matched evaluation and controlled analyses that isolate
    gauge fitting, subspace construction, refitting, and $H_0$ connectivity-scale
    matching, clarifying
    which components account for the observed gains and the regime in which the
    conclusions apply.
\end{itemize}
